\documentclass{article}
\usepackage[margin=0.5in, landscape]{geometry}
\usepackage{booktabs}
\usepackage{array}
\usepackage{longtable}
\usepackage{xcolor}
\usepackage{colortbl}

\begin{document}

\section*{Comprehensive Comparison of Generation II Nuclear Reactors}

\begin{longtable}{>{\raggedright\arraybackslash}p{3.5cm}*{5}{>{\raggedright\arraybackslash}p{3.2cm}}}
\toprule
\rowcolor{gray!30}
\textbf{Characteristic} & \textbf{PWR} & \textbf{BWR} & \textbf{CANDU} & \textbf{Breeder Reactor} & \textbf{AGR \& RBMK} \\
\midrule
\endfirsthead

\toprule
\rowcolor{gray!30}
\textbf{Characteristic} & \textbf{PWR} & \textbf{BWR} & \textbf{CANDU} & \textbf{Breeder Reactor} & \textbf{AGR \& RBMK} \\
\midrule
\endhead

\midrule
\multicolumn{6}{r}{\textit{Continued on next page}} \\
\endfoot

\bottomrule
\endlastfoot

\multicolumn{6}{l}{\textbf{\Large Basic Design Parameters}} \\
\midrule

\textbf{Full Name} & Pressurized Water Reactor & Boiling Water Reactor & CANada Deuterium Uranium Reactor & Fast Breeder Reactor (Liquid Metal Fast Breeder Reactor) & Advanced Gas-cooled Reactor (AGR) \& Reaktor Bolshoy Moshchnosti Kanalny (RBMK) \\
\midrule

\textbf{Development Period} & 1950s-1970s (Gen II: 1965-1996) & 1950s-1970s (Gen II: 1969-1996) & 1960s-1970s & 1950s-1990s & 1960s-1970s (AGR); 1970s-1980s (RBMK) \\
\midrule

\textbf{Moderator Material} & Light water (H$_2$O) - slows neutrons to thermal energies & Light water (H$_2$O) - slows neutrons to thermal energies & Heavy water (D$_2$O) - superior neutron economy & None - operates on fast neutron spectrum & Graphite blocks - allows positive void coefficient (both types) \\
\midrule

\textbf{Primary Coolant} & Pressurized demineralized light water at 325°C & Boiling light water creating steam at 285°C & Pressurized heavy water at 310°C & Liquid sodium (Na) at 550°C or lead-bismuth alloy & Carbon dioxide gas at 640°C (AGR); Light water at 284°C (RBMK) \\
\midrule

\textbf{Fuel Type} & Uranium dioxide (UO$_2$) pellets enriched to 3-5\% U-235 & Uranium dioxide (UO$_2$) pellets enriched to 3-5\% U-235 & Natural uranium dioxide (UO$_2$) or slightly enriched (0.7-0.9\% U-235) & Mixed oxide (MOX) fuel: Pu-239 core with U-238 breeding blanket & Uranium dioxide enriched to 2-3\% (AGR); 2\% enrichment (RBMK) \\
\midrule

\textbf{Fuel Assembly} & Square lattice arrays in pressure vessel, 17×17 or 15×15 rod arrays & Similar to PWR but 8×8 or 10×10 arrays, water gaps for moderation & Horizontal pressure tubes with 37-rod bundles & Hexagonal or square assemblies in sodium pool or loop & Fuel stringers in graphite channels (AGR); Vertical pressure tubes (RBMK) \\
\midrule

\textbf{Core Configuration} & Cylindrical core in single large pressure vessel & Cylindrical core in single large pressure vessel & Horizontal or vertical pressure tubes in calandria vessel & Core in large pool or pipe loops with breeding blanket & Vertical channels in graphite moderator core \\
\midrule

\multicolumn{6}{l}{\textbf{\Large Operational Parameters}} \\
\midrule

\textbf{Operating Pressure} & Primary: 155 bar (2250 psi); Secondary: 70 bar & 75 bar (1085 psi) in reactor vessel & Primary: 100 bar; Moderator: 1-2 bar & Approximately 1 bar for sodium coolant & 40 bar CO$_2$ (AGR); 70 bar water (RBMK) \\
\midrule

\textbf{Operating Temperature} & Primary: 325°C inlet, 330°C outlet; Steam: 280°C & Reactor: 285°C; Steam: 285°C & Primary: 266°C inlet, 310°C outlet; Steam: 260°C & Sodium: 400°C inlet, 550°C outlet; High temp enables efficiency & CO$_2$: 640°C outlet (AGR); Water: 284°C (RBMK) \\
\midrule

\textbf{Steam Generation} & Secondary loop via steam generators (U-tube heat exchangers) & Direct cycle - steam generated in core, goes directly to turbine & Secondary loop with steam generators similar to PWR & Intermediate sodium loop, then water/steam loop (3 loops total) & Secondary loop via heat exchangers (AGR); Direct cycle (RBMK) \\
\midrule

\textbf{Thermal Efficiency} & 32-34\% (limited by saturation temp) & 32-34\% (similar to PWR) & 28-30\% (lower due to heavy water temps) & 38-42\% (higher due to sodium temps) & 41-42\% (AGR - highest of thermal reactors); 31\% (RBMK) \\
\midrule

\textbf{Power Density} & 100-105 kW/liter (high, compact core) & 50-55 kW/liter (lower than PWR) & 10-12 kW/liter (very low, large core) & 200-400 kW/liter (very high) & 2-3 kW/liter (AGR); 3.5 kW/liter (RBMK) - both very low \\
\midrule

\textbf{Refueling Method} & Offline - reactor shutdown, vessel head removed, 18-24 month cycle & Offline - similar to PWR, 18-24 month cycle & Online - fuel channels refueled during operation without shutdown & Offline - complex due to sodium handling, 12-24 month cycle & Online refueling for both AGR and RBMK designs \\
\midrule

\textbf{Fuel Burnup} & 40-50 GWd/tU (gigawatt-days per ton uranium) & 40-50 GWd/tU & 7-8 GWd/tU (lower due to natural U) & 100-200 GWd/tU (very high) & 18-20 GWd/tU (AGR); 20 GWd/tU (RBMK) \\
\midrule

\textbf{Typical Capacity} & 900-1600 MWe (modern designs up to 1700 MWe) & 600-1400 MWe & 540-900 MWe (CANDU-6: 700 MWe) & 250-1200 MWe (BN-800: 800 MWe) & 555-660 MWe (AGR); 1000 MWe (RBMK-1000) \\
\midrule

\textbf{Core Lifetime} & 40-60 years with vessel surveillance & 40-60 years & 30-40 years (pressure tube replacement possible) & 30-40 years (sodium corrosion issues) & 30-40 years (graphite moderation limits) \\
\midrule

\multicolumn{6}{l}{\textbf{\Large Nuclear Physics \& Safety}} \\
\midrule

\textbf{Neutron Spectrum} & Thermal neutrons (0.025 eV) & Thermal neutrons (0.025 eV) & Thermal neutrons with superior moderation & Fast neutrons (100 keV - 1 MeV) - no moderation & Thermal neutrons (graphite moderated) \\
\midrule

\textbf{Neutron Economy} & Moderate - requires enrichment due to absorption in H$_2$O & Moderate - similar to PWR & Excellent - D$_2$O low absorption allows natural U & Breeding ratio: 1.15-1.4 (produces more fuel than consumes) & Good (graphite) but limited by other factors \\
\midrule

\textbf{Reactivity Coefficients} & Negative temperature coefficient; Negative void coefficient (inherently safe) & Negative temperature coefficient; Negative void coefficient & Negative temperature coefficient; Negative void coefficient; Strong negative power coefficient & Small negative or positive depending on design; Delayed neutron fraction important & Negative temp coef. (AGR); POSITIVE void coef. (RBMK - dangerous) \\
\midrule

\textbf{Control Mechanisms} & Control rods (Ag-In-Cd or B$_4$C); Soluble boron in coolant; Chemical shim & Control rods (B$_4$C); Recirculation flow control & Control rods; Adjuster rods; Liquid zone control; Moderator level & Control rods; Reactivity controlled by fuel loading and geometry & Control rods; Some gas flow control \\
\midrule

\textbf{Inherent Safety Features} & Negative feedback coefficients prevent runaway; Double containment; Multiple barriers & Negative void coefficient; Containment structure; ECCS & Strong negative coefficients; Separate moderator and coolant; Low pressure moderator & Double containment for sodium; Passive decay heat removal; Guard vessels & Containment (AGR); RBMK had minimal containment - critical flaw \\
\midrule

\textbf{Emergency Cooling} & High-pressure injection; Accumulators; Low-pressure injection; Multiple redundant systems & High/low pressure ECCS; Automatic depressurization; Core spray systems & Emergency core cooling injection; Moderator cooling; Multiple systems & Decay heat removal via natural circulation in sodium pool & Emergency cooling systems present but varied by design \\
\midrule

\textbf{Containment} & Large dry or ice condenser containment (steel-lined concrete); Full containment & Mark I, II, or III containment designs; Pressure suppression pools & Vacuum building or containment; Lower pressure design & Double containment for sodium and radioactive materials & Concrete biological shield (AGR); RBMK had only partial containment \\
\midrule

\textbf{Void Coefficient} & Strongly negative (loss of coolant reduces reactivity) & Negative (inherently safe) & Strongly negative & Near zero or slightly positive (compensated by Doppler) & Negative (AGR); POSITIVE (RBMK) - led to Chernobyl disaster \\
\midrule

\textbf{Major Accident Risks} & Loss of coolant accident (LOCA); Steam generator tube rupture & LOCA; Recirculation system failure & Loss of coolant/moderator; Pressure tube failure & Sodium fires; Sodium-water reactions; Positive reactivity insertion & Graphite fire (AGR); Power excursion due to positive void (RBMK) \\
\midrule

\multicolumn{6}{l}{\textbf{\Large Advantages \& Disadvantages}} \\
\midrule

\textbf{Key Advantages} & Simple, proven design; Compact core; High power density; Extensive operating experience; Negative safety coefficients; Established supply chain & No steam generators needed; Simpler primary loop; Direct steam cycle more efficient than expected; Lower operating pressure than PWR & Uses natural or low-enriched uranium; Online refueling increases availability; Excellent neutron economy; Can use various fuels (Th, Pu); Fuel flexibility & Breeds more fissile material than consumed; Best uranium utilization (60-80\% vs 1\%); High thermal efficiency; Reduced waste; Transmutes actinides & High thermal efficiency due to gas coolant (AGR); Online refueling capability (both); Continuous operation \\
\midrule

\textbf{Key Disadvantages} & Very high pressure system; Cannot refuel online; Expensive steam generators; Large pressure vessel manufacturing challenge; Boron control complexity & Radioactive steam contaminates turbine; More biological shielding required; Larger containment needed; Limited steam pressure/temp & Physically large size; Heavy water expensive and requires production; Complex piping and maintenance; Lower power density; Tritium production in heavy water & Sodium highly reactive with air and water; Complex handling; Limited commercial success; High capital cost; Fuel reprocessing needed; Proliferation concerns & Large physical size; Complex design (AGR); Positive void coefficient creates instability (RBMK); Graphite fire risk \\
\midrule

\textbf{Economic Factors} & Lower capital cost per MW; Proven supply chain; Standardized designs reduce cost; Operating costs moderate & Similar to PWR; Slightly higher maintenance due to turbine contamination & Higher capital cost; Heavy water cost; Lower capacity factor historically & Very high capital cost; Sodium handling expensive; Limited commercial viability; Reprocessing infrastructure needed & High capital cost; Complex construction; Limited standardization \\
\midrule

\textbf{Maintenance} & Steam generator tube inspection critical; Pressure vessel surveillance; Standard 18-24 month outages & Turbine maintenance complicated by contamination; Recirculation pump maintenance; Jet pump issues & Pressure tube inspection/replacement; Complex coolant systems; Tritium management & Sodium systems require specialized training; Limited repair capability; Component inspection difficult & Complex maintenance; Graphite monitoring; Pressure tube inspection (RBMK) \\
\midrule

\multicolumn{6}{l}{\textbf{\Large Deployment \& Operational History}} \\
\midrule

\textbf{Global Deployment} & ~300 units worldwide (~60\% of all reactors); Most popular design globally & ~80 units worldwide (~20\% of reactors); Second most common & ~50 units worldwide (primarily in Canada, India, China, South Korea, Romania, Argentina) & ~20 built worldwide; Only ~5-8 currently operating; Limited deployment & 14 AGRs in UK (7 still operating); 17 RBMK units built (11 remaining in Russia) \\
\midrule

\textbf{Primary Users} & USA (most reactors), France (58 units), Japan, South Korea, China, Russia, Belgium, Spain & USA, Japan, Sweden, Germany, Switzerland, Finland, Taiwan, Mexico & Canada (inventor), India (largest fleet), China, South Korea, Romania, Argentina & France (Superphénix), Russia (BN-600, BN-800), Japan (Monju - closed), India (FBTR, PFBR) & United Kingdom only (AGR); Russia and Lithuania only (RBMK) \\
\midrule

\textbf{Notable Examples} & Biblis (Germany); Tihange (Belgium); Beaver Valley (USA); Gravelines (France); Most US commercial reactors & Fukushima Daiichi (Japan); Nine Mile Point (USA); Leibstadt (Switzerland); Forsmark (Sweden) & Pickering, Bruce, Darlington (Canada); Rajasthan, Kakrapar (India); Qinshan (China); Cernavodă (Romania); Embalse (Argentina) & Superphénix (France, 1200 MWe, closed); BN-600, BN-800 (Russia, operating); Monju (Japan, closed); Fast Flux Test Facility (USA, closed) & Hinkley Point B, Hartlepool, Heysham, Torness (UK); Chernobyl, Kursk, Leningrad, Smolensk (Russia/former USSR) \\
\midrule

\textbf{Major Incidents} & Three Mile Island (1979, USA) - partial meltdown, contained; Davis-Besse (2002) - vessel head corrosion & Fukushima Daiichi (2011, Japan) - tsunami caused meltdown; Browns Ferry (1975, USA) - fire & No major accidents; Excellent safety record; Minor coolant leaks only & Superphénix multiple sodium leaks; BN-350 minor incidents; Generally safe operation with proper handling & Chernobyl (1986) - catastrophic RBMK explosion and fire; Worst nuclear accident in history \\
\midrule

\textbf{Current Status} & Still the most built and operated type; New Gen III+ PWRs under construction (AP1000, EPR); Continues to dominate new builds & Few new BWRs being built; ABWR (Gen III) operating in Japan; ESBWR (Gen III+) designed but limited orders & Operating successfully; Gen III+ CANDU designs (EC6) being developed; India expanding CANDU fleet; China operating units & Most commercial breeders shut down due to economics; Russia continues development (BN-1200); India developing PFBR; Gen IV concepts include fast reactors & All AGRs scheduled for shutdown by 2030s; All RBMK units to be decommissioned; Chernobyl units 1-3 shut down 1996-2000; No new builds \\
\midrule

\textbf{Future Prospects} & Continues as dominant reactor type; Gen III+ versions being deployed; SMR versions under development & Limited new builds; Gen III ABWR has some success; Market share declining & Stable niche market; Advantages in countries without enrichment; CANDU technology evolving; India committed to expansion & Limited commercial future for current designs; Gen IV fast reactor concepts may revive interest; Integral Fast Reactor (IFR) concept promising & No future for either design; Being replaced by modern reactors; Important historical lessons for safety culture \\
\midrule

\textbf{Lessons Learned} & TMI improved safety culture; Defense-in-depth validated; Importance of operator training; Severe accident management & Fukushima emphasized external hazards; Need for passive safety; Station blackout concerns; Tsunami protection & Excellent safety record demonstrates sound design; Heavy water technology viable; Fuel flexibility valuable & Sodium technology challenging; Economic viability difficult; Breeding not necessary with abundant uranium; Advanced designs needed & RBMK demonstrated dangers of positive void coefficient; Containment essential; Safety culture paramount; AGR showed gas cooling viable but complex \\
\midrule

\multicolumn{6}{l}{\textbf{\Large Technical Specifications}} \\
\midrule

\textbf{Pressure Vessel} & Large forged steel vessel, 4-5m diameter, 12m height, 20-30cm thick, limits plant size & Similar to PWR but different internals, 6m diameter, 22m height & Calandria tank with hundreds of pressure tubes, each 6-10m long & Pool type: large vessel with sodium; Loop type: piping system & Concrete pressure vessel with steel liner (AGR); Graphite stack with channels (RBMK) \\
\midrule

\textbf{Steam Generator} & 2-4 vertical U-tube units, each 20m tall, 50,000-15,000 tubes per unit & None - direct cycle & 4-12 units depending on size & Intermediate heat exchangers (Na-Na, Na-H$_2$O) & Gas-to-water heat exchangers (AGR); None for RBMK \\
\midrule

\textbf{Turbine} & Tandem compound, 4-6 low-pressure units, 1800 rpm (50 Hz) or 3600 rpm (60 Hz) & Similar to PWR but handles radioactive steam & Similar to PWR & Standard turbine with non-radioactive steam & Similar to PWR (AGR); Handles radioactive steam (RBMK) \\
\midrule

\textbf{Containment Volume} & 50,000-75,000 m³ for large dry containment & 50,000-300,000 m³ (Mark III largest) & 70,000-100,000 m³ typical & Double containment for sodium systems & Varies (AGR); Minimal - major design flaw (RBMK) \\
\midrule

\textbf{Spent Fuel Storage} & On-site pool storage for 10-30 years & On-site pool storage & Both wet and dry storage; Can be reused for natural U & Designed for reprocessing; On-site storage limited & On-site storage facilities \\
\midrule

\textbf{Water Consumption} & 90-150 million liters/day for cooling towers or once-through cooling & Similar to PWR & Similar to PWR & Similar to PWR & Lower water consumption due to gas cooling (AGR); Similar to PWR (RBMK) \\
\midrule

\textbf{Site Requirements} & Large cooling water source; Seismic considerations; Population density limits & Similar to PWR & Larger site due to physical size; Heavy water production/storage & Isolated site preferred due to sodium; Seismic isolation important & Large site for complex systems \\
\midrule

\textbf{Construction Time} & 5-7 years typical (Gen II); Modern builds 7-12 years with delays & Similar to PWR & 6-10 years depending on supply chain & 8-15 years (complex construction) & 8-12 years (both types) \\
\midrule

\multicolumn{6}{l}{\textbf{\Large Fuel Cycle \& Waste}} \\
\midrule

\textbf{Uranium Utilization} & ~1\% of natural uranium energy extracted (once-through cycle) & ~1\% of natural uranium (same as PWR) & ~1\% but uses natural uranium (no enrichment needed) & 60-80\% of uranium energy extracted through breeding & ~1-2\% (similar to other thermal reactors) \\
\midrule

\textbf{Plutonium Production} & ~200-300 kg/GWe-year of reactor-grade Pu & Similar to PWR & ~150-200 kg/GWe-year & Breeds weapons-grade Pu-239 from U-238 blanket & Similar to other thermal reactors \\
\midrule

\textbf{Spent Fuel Composition} & 94\% U-238, 1\% U-235, 1\% Pu, 4\% fission products & Similar to PWR & Higher U-235 depletion due to natural U start & Higher Pu content; Lower actinide waste & Similar to PWR and BWR \\
\midrule

\textbf{Reprocessing} & Not required for operation; PUREX available if desired & Not required; Same as PWR & Not required but possible; Advantages for Pu recycling & Required for breeding cycle; PUREX or pyroprocessing & Not typically reprocessed \\
\midrule

\textbf{Waste Volume} & 20-30 m³ spent fuel per GWe-year & Similar to PWR & 90-100 m³ per GWe-year (larger volume due to natural U) & Lower high-level waste volume; Transmutes long-lived actinides & 20-30 m³ per GWe-year \\
\midrule

\textbf{Waste Radiotoxicity} & High-level waste for 10,000+ years & Same as PWR & Similar long-term radiotoxicity & Potentially reduced by actinide transmutation & Similar to PWR and BWR \\
\midrule

\textbf{Proliferation Risk} & Low - reactor-grade Pu difficult for weapons; Safeguarded & Low - same as PWR & Low-moderate - produces Pu but safeguards applied & High - produces weapons-grade Pu-239; Strict controls needed & Low to moderate \\

\bottomrule
\end{longtable}

\section*{Summary Notes}

\textbf{PWR (Pressurized Water Reactor):} The workhorse of the nuclear industry, PWRs represent the most successful commercial reactor design. Their simplicity, compactness, and proven safety record have made them the dominant choice worldwide. The high-pressure operation prevents boiling and allows for compact cores.

\textbf{BWR (Boiling Water Reactor):} The second most common design, BWRs simplify the primary system by allowing boiling in the core, eliminating steam generators. However, this creates radioactive steam that contaminates the turbine, requiring additional shielding.

\textbf{CANDU (Canada Deuterium Uranium):} A unique design that uses heavy water moderation to achieve such excellent neutron economy that natural uranium can be used as fuel. The pressure tube design allows online refueling, maximizing capacity factor. The large size and heavy water costs are offset by fuel cost savings.

\textbf{Breeder Reactor:} Designed to produce more fissile material than consumed, breeders could theoretically extract 60-80\% of uranium's energy versus ~1\% for thermal reactors. However, liquid sodium cooling creates significant engineering challenges and safety concerns. Economic viability has proven elusive, though Russia and India continue development.

\textbf{AGR \& RBMK:} Both use graphite moderation but are otherwise quite different. AGRs use CO$_2$ gas cooling to achieve high thermal efficiency but have complex designs. RBMKs have a fundamentally flawed design with a positive void coefficient that led to the Chernobyl disaster. All RBMKs are being phased out, and AGRs are also being retired in favor of safer designs.


\newpage


\begin{table}[htbp]
\centering
\footnotesize
\begin{tabular}{>{\raggedright\arraybackslash}p{3cm}*{5}{>{\raggedright\arraybackslash}p{3cm}}}
\toprule
\rowcolor{gray!30}
\textbf{Parameter} & \textbf{PWR} & \textbf{BWR} & \textbf{CANDU} & \textbf{Breeder} & \textbf{AGR \& RBMK} \\
\midrule

\multicolumn{6}{l}{\cellcolor{gray!15}\textbf{Core Design}} \\
\textbf{Moderator} & Light water & Light water & Heavy water & None (fast) & Graphite \\
\textbf{Coolant} & Pressurized H$_2$O & Boiling H$_2$O & Heavy water & Liquid Na & CO$_2$ (AGR); H$_2$O (RBMK) \\
\textbf{Fuel} & Enriched U (3-5\%) & Enriched U (3-5\%) & Natural U (0.7\%) & Pu-239 + U-238 & Enriched U (2-3\%) \\
\midrule

\multicolumn{6}{l}{\cellcolor{gray!15}\textbf{Operating Conditions}} \\
\textbf{Pressure} & 155 bar & 75 bar & 100 bar & 1 bar & 40 bar (AGR); 70 bar (RBMK) \\
\textbf{Temperature} & 325°C & 285°C & 310°C & 550°C & 640°C (AGR); 284°C (RBMK) \\
\textbf{Steam Cycle} & Secondary & Direct & Secondary & Tertiary & Secondary (AGR); Direct (RBMK) \\
\textbf{Efficiency} & 33\% & 33\% & 29\% & 40\% & 41\% (AGR); 31\% (RBMK) \\
\textbf{Power Density} & 100 kW/L & 55 kW/L & 12 kW/L & 300 kW/L & 2.5 kW/L \\
\textbf{Refueling} & Offline (18-24 mo) & Offline (18-24 mo) & Online & Offline (12-24 mo) & Online \\
\midrule

\multicolumn{6}{l}{\cellcolor{gray!15}\textbf{Safety \& Physics}} \\
\textbf{Void Coefficient} & Negative & Negative & Negative & Near zero & Negative (AGR); \textcolor{red}{\textbf{Positive (RBMK)}} \\
\textbf{Inherent Safety} & Excellent & Excellent & Excellent & Good & Good (AGR); \textcolor{red}{\textbf{Poor (RBMK)}} \\
\textbf{Containment} & Full double & Full & Full/vacuum & Double (Na) & Full (AGR); \textcolor{red}{\textbf{Partial (RBMK)}} \\
\textbf{Major Accident} & TMI-2 (1979) & Fukushima (2011) & None & Sodium leaks & \textcolor{red}{\textbf{Chernobyl (1986)}} \\
\midrule

\multicolumn{6}{l}{\cellcolor{gray!15}\textbf{Performance}} \\
\textbf{Capacity} & 1000-1600 MWe & 800-1300 MWe & 600-900 MWe & 250-1200 MWe & 500-1000 MWe \\
\textbf{U Utilization} & 1\% & 1\% & 1\% (no enrichment) & 60-80\% & 1-2\% \\
\textbf{Core Lifetime} & 40-60 years & 40-60 years & 30-40 years & 30-40 years & 30-40 years \\
\midrule

\multicolumn{6}{l}{\cellcolor{gray!15}\textbf{Key Advantages}} \\
 & • Compact, proven\\• High power density\\• Strong safety record & • No steam generators\\• Simpler primary loop\\• Direct cycle & • Natural uranium fuel\\• Online refueling\\• Fuel flexibility & • Breeds fissile fuel\\• Best U utilization\\• High efficiency & • High efficiency (AGR)\\• Online refueling\\• Low power density \\
\midrule

\multicolumn{6}{l}{\cellcolor{gray!15}\textbf{Key Disadvantages}} \\
 & • High pressure\\• No online refuel\\• Expensive steam gen. & • Radioactive turbine\\• More shielding\\• Larger containment & • Large physical size\\• Expensive D$_2$O\\• Lower power density & • Reactive Na coolant\\• Complex handling\\• High capital cost & • Complex design (AGR)\\• \textcolor{red}{\textbf{Unstable (RBMK)}}\\• Graphite fire risk \\
\midrule

\multicolumn{6}{l}{\cellcolor{gray!15}\textbf{Deployment}} \\
\textbf{Global Units} & ~300 (60\%) & ~80 (20\%) & ~50 & ~8 operating & 7 AGR; 11 RBMK \\
\textbf{Primary Users} & USA, France, China, Japan & USA, Japan, Sweden & Canada, India, China & Russia, India & UK; Russia \\
\textbf{Status} & Dominant, expanding & Declining & Stable niche & Limited future & Phase out \\

\bottomrule
\end{tabular}
\caption{Comprehensive Comparison of Generation II Nuclear Reactors}
\end{table}

\section*{Key Takeaways}

\textbf{PWR:} Most successful design (60\% of global fleet). Compact, safe, proven technology. Standard for new builds.

\textbf{BWR:} Second most common (20\%). Simpler than PWR but radioactive steam is a drawback. Fukushima highlighted vulnerability to external hazards.

\textbf{CANDU:} Unique heavy water design uses natural uranium—no enrichment needed. Excellent for countries without enrichment capability. Online refueling maximizes availability.

\textbf{Breeder:} Extracts 60-80× more energy from uranium but liquid sodium creates operational challenges. Limited commercial success despite potential benefits.

\textbf{AGR \& RBMK:} AGR achieved high efficiency with gas cooling but complex. RBMK's positive void coefficient caused Chernobyl—worst nuclear accident in history. Both being phased out.

\end{document}

